% Compilation : pdflatex -shell-escape dm.tex
\documentclass[11pt,a4paper]{article}

\usepackage[utf8]{inputenc}
\usepackage[T1]{fontenc}
\usepackage[french]{babel}
\usepackage[top=2.5cm, bottom=2.5cm, left=2.8cm, right=2.8cm]{geometry}
\usepackage{minted}
\usepackage{xcolor}
\usepackage[framemethod=TikZ]{mdframed}
\usepackage{amsmath}
\usepackage{amssymb}
\usepackage[hidelinks]{hyperref}
\usepackage{parskip}
\usepackage{titlesec}
\usepackage{enumitem}
\usepackage{microtype}
\usepackage{fontawesome5}

% ── Couleurs ──────────────────────────────────────────────────────────────────
\definecolor{partcolor}{RGB}{25, 65, 140}
\definecolor{codebg}{RGB}{250, 250, 252}
\definecolor{rappelbg}{RGB}{255, 252, 230}
\definecolor{rappelframe}{RGB}{190, 155, 40}
\definecolor{exobg}{RGB}{240, 247, 255}
\definecolor{exoframe}{RGB}{55, 115, 205}
\definecolor{bonusbg}{RGB}{242, 252, 242}
\definecolor{bonusframe}{RGB}{60, 160, 80}
\definecolor{todofg}{RGB}{175, 50, 50}

% ── Titres de sections ─────────────────────────────────────────────────────────
\titleformat{\section}
{\large\bfseries\color{partcolor}}
{\thesection.}{0.5em}{}[\vspace{-4pt}\rule{\linewidth}{1.2pt}]

\titleformat{\subsection}
{\normalsize\bfseries}
{}{}{}

% ── Boîtes ────────────────────────────────────────────────────────────────────
\mdfdefinestyle{rappelstyle}{
  backgroundcolor=rappelbg,
  linecolor=rappelframe,
  linewidth=1.2pt,
  roundcorner=4pt,
  innerleftmargin=10pt,
  innerrightmargin=10pt,
  innertopmargin=6pt,
  innerbottommargin=6pt,
  skipabove=8pt,
  skipbelow=8pt,
}

\mdfdefinestyle{exostyle}{
  backgroundcolor=exobg,
  linecolor=exoframe,
  linewidth=2pt,
  topline=false,
  bottomline=false,
  rightline=false,
  innerleftmargin=12pt,
  innerrightmargin=8pt,
  innertopmargin=6pt,
  innerbottommargin=6pt,
  skipabove=6pt,
  skipbelow=6pt,
}

\mdfdefinestyle{bonusstyle}{
  backgroundcolor=bonusbg,
  linecolor=bonusframe,
  linewidth=2pt,
  topline=false,
  bottomline=false,
  rightline=false,
  innerleftmargin=12pt,
  innerrightmargin=8pt,
  innertopmargin=6pt,
  innerbottommargin=6pt,
  skipabove=6pt,
  skipbelow=6pt,
}

\newenvironment{rappel}{
\begin{mdframed}[style=rappelstyle]}{
\end{mdframed}}
\newenvironment{exo}{
\begin{mdframed}[style=exostyle]}{
\end{mdframed}}
\newenvironment{bonus}{
\begin{mdframed}[style=bonusstyle]}{
\end{mdframed}}

% ── Code Coq ──────────────────────────────────────────────────────────────────
\setminted[coq]{
  bgcolor=codebg,
  fontsize=\small,
  breaklines=true,
  breakanywhere=false,
  autogobble=true,
}

% Commande pour code inline
\newcommand{\cq}[1]{\mintinline{coq}{#1}}

% ── Page de titre ──────────────────────────────────────────────────────────────
\title{%
  \vspace{-1.5cm}%
  \textbf{Devoir Maison : Inverser deux fois une liste}\\[0.4em]
  \large Logique et Fondements de l'Informatique%
}
\date{}
\author{}

% ══════════════════════════════════════════════════════════════════════════════
\begin{document}
\maketitle
\vspace{-1.2cm}

\begin{rappel}
  \textbf{Modalités de rendu.}

  Le fichier à compléter est \texttt{dm.v}. C'est un fichier texte contenant du
  code Rocq. Voici la marche à suivre :

  \begin{enumerate}[noitemsep,topsep=4pt]
    \item \textbf{Ouvrez jsCoq} dans votre navigateur
      (\url{https://jscoq.github.io/scratchpad.html}).
    \item \textbf{Copiez-collez} le contenu du fichier \texttt{dm.v} dans
      l'éditeur jsCoq.
    \item \textbf{Travaillez dans jsCoq} : complétez les preuves et les réponses
      aux questions directement dans l'éditeur. Vous pouvez exécuter le code pas
      à pas pour vérifier vos preuves au fur et à mesure.
    \item \textbf{Enregistrez} votre travail en cliquant sur le bouton
      de sauvegarde (\faDownload) dans jsCoq. Cela téléchargera un fichier
      \texttt{.v} sur votre ordinateur.
    \item \textbf{Rendez} le fichier \texttt{.v} téléchargé par email à
      \href{mailto:pablo.donato@tuta.io}{\texttt{pablo.donato@tuta.io}}.
  \end{enumerate}

  Tout votre travail (preuves \emph{et} réponses aux questions de compréhension)
  doit se trouver dans ce fichier unique. Les réponses en français se rédigent
  dans des commentaires Coq, entre \cq{(*} et \cq{*)}.
\end{rappel}

\bigskip

Ce DM vous guide vers la preuve d'un résultat classique de vérification
formelle :
\[
  \forall\, A,\; \forall\, l : \mathrm{list}\;A, \quad
  \cq{rev}(\cq{rev}\; l) = l
\]
c'est-à-dire qu'inverser deux fois une liste redonne la liste de départ.

Cette preuve ne peut pas se faire directement : elle nécessite des résultats intermédiaires, appelés
\textbf{lemmes auxiliaires}. En mathématiques comme en programmation, on décompose un problème
complexe en sous-problèmes plus simples. En Rocq, cela se traduit par des lemmes qu'on prouve
séparément, puis qu'on réutilise avec la tactique \cq{rewrite}\footnote{On utilise plus généralement
  la tactique \cq{apply} si la conclusion du lemme n'est pas une égalité. Dans ce DM on prouve
uniquement des égalités.}.

Les exercices 1 à 3 établissent les lemmes dont vous aurez besoin pour
l'exercice~4. Remplacez \cq{admit} par votre preuve. Quand vous avez terminé
un exercice, remplacez \cq{Admitted} par \cq{Qed} --- Rocq vérifiera que
votre preuve est complète.

\begin{rappel}
  \textbf{Rappel des tactiques utiles :}
  \begin{itemize}[noitemsep,topsep=2pt]
    \item \cq{intros} : introduire les hypothèses
    \item \cq{induction l as [| h t IH]} : raisonner par induction sur une liste
    \item \cq{simpl} : simplifier (appliquer les définitions)
    \item \cq{rewrite lemme} : réécrire le but en utilisant un lemme ou une hypothèse
    \item \cq{rewrite <- lemme} : réécrire de droite à gauche
    \item \cq{reflexivity} : conclure quand les deux côtés sont égaux
  \end{itemize}

\end{rappel}

\begin{rappel}
  \textbf{Définitions de référence.} On utilise \cq{rev} et \cq{app} (noté \cq{++})
  de la bibliothèque standard :
\begin{minted}{coq}
  Fixpoint app {A : Type} (l1 l2 : list A) : list A :=
    match l1 with
    | []     => l2
    | h :: t => h :: (app t l2)
    end.

  Fixpoint rev {A : Type} (l : list A) : list A :=
    match l with
    | []     => []
    | h :: t => rev t ++ [h]
    end.
\end{minted}
  %   Le préambule du fichier \texttt{dm.v} est :
% \begin{minted}{coq}
%   Require Import List.
%   Import ListNotations.
% \end{minted}
\end{rappel}

\begin{rappel}
  \textbf{Exemple : utiliser \cq{rewrite} avec un lemme.}

  Supposons qu'on ait déjà prouvé le lemme suivant :
\begin{minted}{coq}
  Lemma add_0_r : forall n : nat, n + 0 = n.
\end{minted}
  Dans une preuve ultérieure, si le but contient \cq{n + 0}, on peut utiliser
  \cq{rewrite add_0_r} pour le remplacer par \cq{n} :
\begin{minted}{coq}
  (* But courant :  f (n + 0) = f n *)
  rewrite add_0_r.
  (* Nouveau but :  f n = f n *)
  reflexivity.
\end{minted}
  De même, \cq{rewrite <- add_0_r} ferait le remplacement inverse :
  de \cq{n} vers \cq{n + 0}.

  % Dans ce DM, les exercices 1 à 3 établissent des lemmes que vous
  % réutiliserez de cette manière dans les exercices suivants.
\end{rappel}

% ══════════════════════════════════════════════════════════════════════════════
\section{Concaténer avec la liste vide ne fait rien}

% \textbf{Énoncé :}
% $\forall\, A,\; \forall\, l : \mathrm{list}\;A, \quad l \mathbin{\texttt{++}} [\,] = l$

% \medskip

\textit{Indice :} \cq{app} est défini par récursion sur son \emph{premier}
argument. Donc \cq{[] ++ l} se réduit par définition, mais \cq{l ++ []} ne se
réduit pas si \cq{l} est une variable. Procédez par induction sur \cq{l}.

\begin{exo}
\begin{minted}{coq}
  Lemma app_nil_r : forall (A : Type) (l : list A),
      l ++ [] = l.
\end{minted}
\end{exo}

% ══════════════════════════════════════════════════════════════════════════════
\section{La concaténation est associative}

% \textbf{Énoncé :}
% $\forall\, A,\; \forall\, l_1, l_2, l_3 : \mathrm{list}\;A, \quad
% (l_1 \mathbin{\texttt{++}} l_2) \mathbin{\texttt{++}} l_3 =
% l_1 \mathbin{\texttt{++}} (l_2 \mathbin{\texttt{++}} l_3)$

% \medskip

\textit{Indice :} même structure que l'exercice 1, par induction sur \cq{l1}.

\begin{exo}
\begin{minted}{coq}
  Lemma app_assoc : forall (A : Type) (l1 l2 l3 : list A),
      (l1 ++ l2) ++ l3 = l1 ++ (l2 ++ l3).
\end{minted}
\end{exo}

% ══════════════════════════════════════════════════════════════════════════════
\section{Inverser une concaténation}

% \textbf{Énoncé :}
% $\forall\, A,\; \forall\, l_1, l_2 : \mathrm{list}\;A, \quad
% \cq{rev}\;(l_1 \mathbin{\texttt{++}} l_2) =
% \cq{rev}\;l_2 \mathbin{\texttt{++}} \cq{rev}\;l_1$

% \medskip

\textit{Indice :} induction sur \cq{l1}. Vous aurez besoin des exercices 1 et 2.

\begin{exo}
\begin{minted}{coq}
  Lemma rev_app : forall (A : Type) (l1 l2 : list A),
      rev (l1 ++ l2) = rev l2 ++ rev l1.
\end{minted}
\end{exo}

% ══════════════════════════════════════════════════════════════════════════════
\section{Inverser deux fois redonne la liste de départ}

% \textbf{Énoncé :}
% $\forall\, A,\; \forall\, l : \mathrm{list}\;A, \quad
% \cq{rev}\;(\cq{rev}\;l) = l$

% \medskip

C'est le résultat principal.

\textit{Indice :} induction sur \cq{l}. Vous aurez besoin de l'exercice 3.

\begin{exo}
\begin{minted}{coq}
  Theorem rev_involutive : forall (A : Type) (l : list A),
      rev (rev l) = l.
\end{minted}
\end{exo}

% ══════════════════════════════════════════════════════════════════════════════
\section*{Questions de compréhension}

Répondez aux questions suivantes en quelques phrases, directement dans le
fichier \texttt{dm.v}, en commentaires Coq (entre \cq{(*} et \cq{*)}).

\begin{bonus}
  \textbf{Question 1.}
  Pourquoi ne peut-on pas prouver \cq{rev (rev l) = l} simplement avec
  \cq{reflexivity}, sans induction ?

  \medskip

  \textbf{Question 2.}
  Quel est le rôle du lemme \cq{rev_app} dans la preuve de \cq{rev_involutive} ?
  Pourquoi est-il nécessaire ?

  \medskip

  \textbf{Question 3.}
  Les exercices 1 à 3 sont des lemmes auxiliaires. Aurait-on pu s'en passer et
  tout prouver d'un seul coup dans l'exercice 4 ? Quel est l'intérêt de
  décomposer la preuve en lemmes ?
\end{bonus}

\end{document}
